\documentclass[11pt,a4paper]{article}
\usepackage[margin=25mm]{geometry}
\usepackage{amsmath,amssymb,siunitx,hyperref,listings,microtype}
\hypersetup{colorlinks=true,linkcolor=blue,urlcolor=blue,pdftitle={Report 7 - Generator Model Family for OpenHPL/OpenHPLjl},pdfauthor={Madhusudhan Pandey}}
\lstset{basicstyle=\ttfamily\small,breaklines=true,frame=single}
\title{Report 7 - Generator Model Family for OpenHPL/OpenHPLjl}
\author{Madhusudhan Pandey}
\date{19 September 2026}
\begin{document}
\maketitle
\begin{abstract}
This report introduces the generator model family. The first implementation is a reduced algebraic mechanical-to-electrical power converter; later models will add swing dynamics, rotor angle, voltage, and synchronous-machine states.
\end{abstract}
\section{Updated Context}
The generator converts shaft mechanical power into electrical active power:
\[
\boxed{\text{Shaft}\rightarrow\text{Generator}\rightarrow\text{Grid}}.
\]

\section{Generalized Concept Model}
Inputs are shaft speed and torque. The principal electrical output is active power:
\[
\omega,\tau\rightarrow\text{Generator}\rightarrow P_e.
\]

\section{Other Models Available in the Literature}
A useful hierarchy is
\[
\boxed{\text{ideal converter}\rightarrow\text{swing model}\rightarrow\text{classical synchronous machine}\rightarrow\text{higher-order dq machine}}.
\]

\section{Mathematics Involved}
Mechanical input power is
\[
P_m=\tau\omega.
\]
The first reduced generator uses
\[
\boxed{P_e=\eta_gP_m}.
\]
Later synchronous-machine models will introduce
\[
M\dot\omega=P_m-P_e-D\Delta\omega
\]
and rotor-angle dynamics.

\section{OpenHPL-Specific Modeling Interpretation}
OpenHPL supports turbine-generator mechanical coupling and power-system integration. OpenHPLjl begins with a transparent reduced converter before adding detailed machine equations.

\section{Implementation in Julia ModelingToolkit / SciML}
\begin{lstlisting}
@component function IdealGenerator(; name, eta=0.98)
    @named shaft = RotationalPort()
    @named grid = ElectricalPowerPort()

    Pm ~ shaft.tau * shaft.omega
    Pe ~ eta * Pm
    grid.P ~ -Pe
    grid.omega ~ shaft.omega
end
\end{lstlisting}

\section{Model Assembly and Connections}
The shaft rotational port connects to the generator mechanical port; the electrical power port connects to an infinite-grid or future dynamic-grid model.

\section{Numerical Experiment / Validation}
For
\[
\tau=1000~\si{\newton\meter},\quad
\omega=100~\si{\radian\per\second},\quad
\eta_g=0.98,
\]
\[
P_m=100~\si{\kilo\watt},
\qquad
\boxed{P_e=98~\si{\kilo\watt}}.
\]

\section{Expected Outputs}
The minimum outputs are
\[
P_m(t),\quad P_e(t),\quad\omega(t).
\]

\section{Engineering Interpretation}
This reduced model is useful for early energy-balance and frequency-control assembly. Detailed synchronous-machine states should only be added when voltage and rotor-electromagnetic dynamics are required.

\section*{Conclusion}
The generator family now has a clean starting interface. Next models are a swing-equation generator and a classical synchronous generator.
\end{document}
